\section{All rational parameters in target height}\label{sec:rational-counting}

The local criterion also permits inert prime cubes in either numerator or
denominator.  Counting the resulting rational parameters requires two
additional ingredients: coprimality of numerator and denominator, and the
content of the projective target coordinates.  The latter changes the
leading constant at both \(2\) and \(3\).

\begin{theorem}\label{thm:rational-count}
Let \(\ell=\{y_a:a\in\Q\}\) be the target line in
\eqref{eq:target-family}, and let
\[
 N_{\Phi,\ell}(B)=
 \#\{a\in\Q:\ H(y_a)\le B,
       \ \Phi^{-1}(y_a)\text{ fails the Hasse principle}\}.
\]
Then
\begin{equation}
 N_{\Phi,\ell}(B)\sim C_\Phi\frac{B^2}{\log B},
 \qquad
 C_\Phi=
 \frac{11}{1728\sqrt3\,\zeta(2)}
 \prod_{p\equiv2\ (3)}\frac{1+p^{-3}}{1-p^{-3}}.
 \label{eq:rational-constant}
\end{equation}
The product is absolutely convergent and positive.
\end{theorem}

\subsection{Primitive coordinates and the height region}

Write \(a=m/n\), where \(m\ne0\), \(n>0\), and
\(\gcd(m,n)=1\).  The projective coordinates are
\[
 [1:y_{m/n}]=[9n:-9n:32m:24m+3n].
\]
Their positive common divisor is exactly
\[
 d(m,n)=\gcd(n,8)\gcd(m,3).
\]
Indeed, a common prime divisor other than \(2,3\) would divide both
\(m\) and \(n\).  At \(3\), the common valuation is one if
\(3\mid m\), and zero otherwise: when \(3\mid m\), the last
coordinate has valuation one because \(3\nmid n\).  At \(2\), there
is no common factor if \(n\) is odd.  If \(n\) is even, then \(m\)
is odd and the common valuation is \(\min\{v_2(n),3\}\), as is seen
from \(9n,32m,24m+3n\).  Consequently
\begin{equation}\label{eq:rational-height}
 H(y_{m/n})=
 \frac{\max\{9n,32|m|,|24m+3n|\}}
      {\gcd(n,8)\gcd(m,3)}.
\end{equation}

Introduce the bounded polygon
\[
 \Omega=\{(x,y)\in\R^2:y>0,
       \ \max(9y,32|x|,|24x+3y|)\le1\}.
\]
Its part with \(x<0\) is a rectangle of area \(1/288\).  Its part
with \(x>0\) is the rectangle
\((0,1/32]\times(0,1/9]\), with the upper right triangle removed by
\(24x+3y\le1\).  The triangle has horizontal and vertical legs
\(1/288\) and \(1/36\), respectively.  Thus
\begin{equation}\label{eq:rational-area}
 \operatorname{area}(\Omega)=
 \frac1{288}+\frac1{288}-\frac1{20736}
 =\frac{143}{20736}.
\end{equation}

\subsection{Coprime pairs and the cubic residue condition}

Define a multiplicative indicator \(g\) on positive integers by
\[
 g(t)=1\quad\Longleftrightarrow\quad
 (t,6)=1\ \text{ and }\quad
 v_p(t)\equiv0\pmod3\ \text{for every }p\equiv2\pmod3.
\]
Its Dirichlet series and leading constant are
\begin{align*}
 F(s)&=
 \prod_{p\equiv1\ (3)}(1-p^{-s})^{-1}
 \prod_{\substack{p\equiv2\ (3)\\p>2}}(1-p^{-3s})^{-1},\\
 K&=\lim_{s\to1^+}\frac{(s-1)^{1/2}F(s)}{\sqrt\pi}
   =\frac{G_3(1)}{\sqrt\pi}
     \prod_{\substack{p\equiv2\ (3)\\p>2}}(1-p^{-3})^{-1}.
\end{align*}
Let \(\chi\) be either character of order three modulo \(9\).
On units, its kernel is \(\{1,-1\}\), and therefore
\[
 \mathbf1_{u/v\equiv\pm1\ (9)}
 =\frac13\sum_{j=0}^2\chi^j(u)\overline{\chi^j(v)}
 \qquad ((uv,3)=1).
\]
For completeness, the required univariate estimates are
\begin{equation}\label{eq:rational-sd}
 \sum_{t\le X}g(t)\sim K\frac{X}{\sqrt{\log X}},
 \qquad
 \sum_{t\le X}g(t)\chi^j(t)
 =o\left(\frac{X}{\sqrt{\log X}}\right)\quad(j=1,2).
\end{equation}
They follow from \cite[Theorem~1.2]{BretecheTenenbaum}, with prime-mean
parameter \(\rho=1/2\) in the first case and \(\rho=0\) in the other
two.  Here all prime-power values have absolute value at most one; the
class \(\mathcal F(r,\sigma)\) in that theorem applies with
\(r=1/2\), \(\sigma=3/4\).  The prime number theorem in the fixed
classes modulo \(9\) gives its prime-sum hypothesis for every fixed
\(A>0\): the values of either nontrivial cubic character sum to zero
on \(\{1,4,7\}\).  The theorem's coefficients vanish for \(\rho=0\).
Its principal leading coefficient is the displayed limit defining
\(K\).  In particular, for all \(X\ge1\),
\begin{equation}\label{eq:rational-sd-upper}
 \sum_{t\le X}g(t)\ll\frac{X}{\sqrt{\log(2X)}}.
\end{equation}

To impose coprimality, put
\[
 \mathcal E(z)=
 \prod_{p\equiv1\ (3)}(1-p^{-z})
 \prod_{\substack{p\equiv2\ (3)\\p>2}}(1-p^{-3z})
 =\sum_{d\ge1}\frac{h(d)}{d^z}.
\]
This series converges absolutely for \(\Re z>1\).  For
\(\eta\in\{1,\chi,\chi^2\}\), Euler-factor comparison gives
\begin{align*}
 &\sum_{\substack{u,v\ge1\\(u,v)=1}}
 \frac{g(u)g(v)\eta(u)\overline{\eta(v)}}{u^s v^t}\\
 &\hspace{15mm}=
 \left(\sum_{u\ge1}\frac{g(u)\eta(u)}{u^s}\right)
 \left(\sum_{v\ge1}\frac{g(v)\overline{\eta(v)}}{v^t}\right)
 \mathcal E(s+t),\qquad \Re s,\Re t>1.
\end{align*}
At a split prime the correction is \(1-p^{-s-t}\); at an odd inert
prime it is \(1-p^{-3s-3t}\).  The character factors cancel on the
diagonal in both cases.

It follows by coefficient comparison that the corresponding pair sum in
\(u\le X,v\le Y\) equals
\[
 \sum_{d\le\min(X,Y)}h(d)
 \left(\sum_{u\le X/d}g(u)\eta(u)\right)
 \left(\sum_{v\le Y/d}g(v)\overline{\eta(v)}\right).
\]
Take \(X=\alpha B,Y=\beta B\), with fixed positive
\(\alpha,\beta\).  For \(d\le\sqrt B\),
\eqref{eq:rational-sd-upper} bounds the summand by
\(O_{\alpha,\beta}(B^2|h(d)|/(d^2\log B))\).
For \(d>\sqrt B\), the trivial bounds give a total
\[
 \ll B^2\sum_{d>\sqrt B}\frac{|h(d)|}{d^2}
 \ll B^{7/4},
\]
using \(\sum |h(d)|d^{-3/2}<\infty\).  Dominated convergence and
\eqref{eq:rational-sd} therefore give the principal main term
\(K^2\mathcal E(2)\alpha\beta B^2/\log B\), while the other two
twists are \(o(B^2/\log B)\).

After applying the character projector, we obtain, for any fixed bounded
Jordan measurable region \(D\subset(0,\infty)^2\),
\begin{equation}\label{eq:rational-polygon-count}
 \#\left\{\begin{array}{@{}l@{}}
 (u,v)\in BD\cap\Z^2:g(u)g(v)=1,\\
 (u,v)=1,\quad u/v\equiv\pm1\pmod9
 \end{array}\right\}
 \sim\frac{K^2\mathcal E(2)}3\operatorname{area}(D)
       \frac{B^2}{\log B},
\end{equation}
provided the area is positive.  Indeed, inclusion--exclusion first gives
this assertion for rectangles with arbitrary fixed positive endpoints.
Upper and lower finite grid approximations then give it for \(D\).
The rectangles adjoining the axes have arbitrarily small total area, and
are controlled by the same rectangle estimate.  This also proves the
corresponding \(o(B^2/\log B)\) assertion when the area is zero.

\subsection{The valuations at two and three}

Every reduced parameter satisfying the local criterion has a unique
expression
\[
 m=\varepsilon\,8^i27^j u,\qquad
 n=8^k27^l v,
 \qquad \varepsilon\in\{1,-1\},\quad i,j,k,l\ge0,
\]
where \(ik=jl=0\), \(g(u)g(v)=1\), and \((u,v)=1\).
The unit condition at \(3\) is exactly
\(u/v\equiv\pm1\pmod9\): powers of \(8\) and the sign are cubes
modulo \(9\), and the powers of \(3\) have been removed.
For this stratum write
\[
 r=8^i27^j,\qquad s=8^k27^l,\qquad
 d=8^{\mathbf1_{k>0}}3^{\mathbf1_{j>0}}.
\]
The height condition is \((\varepsilon r u,s v)\in dB\Omega\).
For fixed indices the two sign regions in \((u,v)\) thus have combined
area \(d^2\operatorname{area}(\Omega)/(rs)\).
Equation~\eqref{eq:rational-polygon-count} applies to each of them.

The passage from finitely many strata to all strata is justified as
follows.  The height bounds imply \(u\le CB/r,v\le CB/s\), with an
absolute constant \(C\), because \(d\le24\).  If
\(r,s\le B^{1/4}\), equation~\eqref{eq:rational-sd-upper}, after
discarding coprimality and the residue condition, bounds the count by
\[
 \ll\frac{B^2}{rs\log B}.
\]
The sum of \(1/(rs)\) over all allowed indices converges.  On the
remaining strata, the trivial bound is \(O(B^2/(rs))\); moreover
\[
 \sum_{\max(r,s)>B^{1/4}}\frac1{rs}
 \le B^{-1/8}\sum_{i,j,k,l\ge0}(8^{i+k}27^{j+l})^{-1/2}
 \ll B^{-1/8}.
\]
These strata contribute \(o(B^2/\log B)\).  Dominated convergence
therefore permits summation of the fixed-stratum main terms.

The resulting valuation weight factors as
\begin{align*}
 \sum_{\substack{i,j,k,l\ge0\\ik=jl=0}}\frac{d^2}{rs}
 &=\left(1+\sum_{i\ge1}8^{-i}
             +64\sum_{k\ge1}8^{-k}\right)
   \left(1+9\sum_{j\ge1}27^{-j}
             +\sum_{l\ge1}27^{-l}\right)\\
 &=\frac{72}{7}\cdot\frac{18}{13}
   =\frac{1296}{91}.
\end{align*}
The factor \(64\) comes from denominator powers of \(2\), and the
factor \(9\) from numerator powers of \(3\); both are squares of
the content corrections in \eqref{eq:rational-height}.

\subsection{Evaluation of the constant and exclusion of cubes}

Using the expression for \(G_3(1)\) in Section~\ref{sec:counting},
\(L(1,\chi_{-3})=\pi/(3\sqrt3)\), and
\(\prod_{p\ne3}(1-p^{-2})=9/(8\zeta(2))\), we find
\begin{align*}
 K^2\mathcal E(2)
 &=\frac1{4\sqrt3\,\zeta(2)}
   \prod_{\substack{p\equiv2\ (3)\\p>2}}
          \frac{1+p^{-3}}{1-p^{-3}}\\
 &=\frac7{36\sqrt3\,\zeta(2)}
   \prod_{p\equiv2\ (3)}\frac{1+p^{-3}}{1-p^{-3}}.
\end{align*}
The \(L\)-value can be evaluated without further analytic input:
\[
 L(1,\chi_{-3})=
 \sum_{k\ge0}\left(\frac1{3k+1}-\frac1{3k+2}\right)
 =\int_0^1\frac{dx}{1+x+x^2}
 =\frac{\pi}{3\sqrt3}.
\]
Combining the factor \(1/3\) for the unit condition, the polygon area,
and the valuation weight, the constant is
\[
 \frac13\cdot\frac{143}{20736}\cdot\frac{1296}{91}
 \cdot\frac7{36\sqrt3\,\zeta(2)}
 \prod_{p\equiv2\ (3)}\frac{1+p^{-3}}{1-p^{-3}}
 =C_\Phi.
\]

This calculation has counted nonzero parameters with a point over every
completion, including global cubes.  By \eqref{eq:rational-height},
\(H(y_{m/n})\le B\) implies \(|m|\le3B/4\) and \(n\le8B/3\).
If \(m/n\) is a rational cube in lowest terms, then
\(m=b^3,n=c^3\) with \((b,c)=1\); there are only \(O(B^{2/3})\)
such pairs.  Removing them leaves exactly the Hasse-failing parameters by
the local classification.  The parameter \(a=0\) does not fail the
Hasse principle because its actual fiber has no real point, as shown in
Theorem~\ref{thm:rational-classification}.  Finally, \(a\mapsto y_a\) is
injective by its second coordinate.  The parameter count is therefore
the target count asserted in Theorem~\ref{thm:rational-count}.
